% !TEX root = ../main.tex

\section{Related Works}
\label{sec:relatedworks}

Personalized federated learning for clients with non-IID data has attracted much attention~\cite{deng2020adaptive, fallah2020personalized,  kulkarni2020survey, mansour2020three}. Particularly, our work is related to global federated learning, local customization and multi-task federated learning.

\nop{
However, it has been empirically shown to diverge on non-IID data across clients~\cite{mcmahan2016communication, li2018federated, zhao2018federated}. 

and Zhao \textit{et~al.}~\cite{zhao2018federated} globally share a small subset of data across all clients, which, however, infringes the data privacy of clients.
}

Global federated learning~\cite{ji2019learning,  mcmahan2016communication, wang2020federated,   yurochkin2019bayesian} trains a single global model to minimize an empirical risk function over the union of the data across all clients.
When the data is non-IID across different clients, however, it is difficult to converge to a good global model that achieves a good personalized performance on every client~\cite{kairouz2019advances, li2018federated, mcmahan2016communication, zhao2018federated}.

%\textbf{Global federated learning.}
\nop{
Global federated learning trains a single global model to minimize the empirical risk function over the union of the data across all clients.
FedAvg~\cite{mcmahan2016communication} produces a global model by simply averaging the parameters of local models.
FedProx~\cite{li2018federated} improves the robustness of FedAvg to non-IID data by keeping local models close to the global model.
FedAtt~\cite{ji2019learning} aggregates local models into a global model by a layer-wise attention mechanism between neural networks. PFNM~\cite{yurochkin2019bayesian} matches the neurons of client neural networks before averaging them, and FedMA~\cite{wang2020federated} extends PFNM to more complicated deep neural networks such as CNNs and LSTMs.
Most global federated learning methods use a single global model for all clients. However, when the data is non-IID across different clients, it is difficult to converge to a good global model that universally achieves a good personalized performance on every client~\cite{kairouz2019advances, li2018federated, mcmahan2016communication, zhao2018federated}.
}

\nop{
\begin{itemize}
    \item As one of the classical global federated learning methods, FedAvg~\cite{mcmahan2016communication} runs stochastic gradient descent to train local models in parallel and averages 
    the parameters of local models only once in a while to produce a global model. 
    FedAvg has been empirically shown to diverge for non-IID data across clients~\cite{mcmahan2016communication, li2018federated, zhao2018federated}, thus FedProx~\cite{li2018federated} improves the robustness of FedAvg to device heterogeneity and non-IID data distribution by adding a proximal term to keep local models close to the global model.
    Zhao \textit{et~al.}~\cite{zhao2018federated} improved the robustness to non-IID data distribution by globally sharing a small subset of data across all clients, which, however, infringes the data privacy of clients.
    
    \item PFNM~\cite{yurochkin2019bayesian} addresses the invariance problem of fully connected feedforward networks by matching the neurons of client neural networks before averaging them. FedMA~\cite{wang2020federated} further extends PFNM to more complicated deep neural networks by exploring the layers-wise invariance of CNNs and LSTMs.
    
    \nop{
    \item Ji \textit{et~al.}~\cite{ji2020dynamic} improves the communication efficiency of federated learning by dynamically sampling collaborating clients and selectively masking neural parameters of client models.
    }
\end{itemize}
}

\nop{
Most horizontal federated learning use one model for all clients. Under non-IID settings, this leads to convergence problems and degenerated prediction performance.

The original goal of federated learning, training a single global model on the union of client data sets, becomes
harder with non-IID data.


For all these methods, using one global model is hard to tackle the statistical challenge brought by non-IID data distributions on different clients.
}

\nop{
While [370, 428, 399, 371] mainly focused on the IID case, the analysis technique can be extended to
the non-IID case by adding an assumption on data dissimilarities, for example by constraining the difference
between client gradients and the global gradient [266, 261, 265, 401] or the difference between client and
global optimum values [264, 232]. Under this assumption, Yu \textit{et~al.} [429] showed that the error bound
of local SGD in the non-IID case becomes worse. In order to achieve the rate of 1=
p
TKN (under nonconvex
objectives), the number of local updates K should be smaller than T1=3=N, instead of T=N3 as in
the IID case [399]. Li \textit{et~al.} [261] proposed to add a proximal term in each local objective function so as
to make the algorithm be more robust to the heterogeneity across local objectives. The proposed FedProx
algorithm empirically improves the performance of federated averaging. However, it is unclear whether it
could provably improve the convergence rate. Khaled \textit{et~al.} [232] assumes all clients participate, and uses
batch gradient descent on clients, which can potentially converge faster than stochastic gradients on clients.
}

%\textbf{Local customization of a global model.}
Local customization methods~\cite{ chen2018federated, fallah2020personalized, jiang2019improving, khodak2019adaptive, kulkarni2020survey, mansour2020three, nichol2018first, schneider2019mass, wang2019federated} build a personalized model for each client by customizing a well-trained global model.  There are several ways to conduct customization. 
A practical way to customize a personalized model is local fine-tuning~\cite{ben2010theory, cortes2014domain, mansour2020three, mansour2009domain,  schneider2019mass, wang2019federated}, where the global model is fine-tuned using the private data of each client to produce a personalized model for the client.
Similarly, meta-learning methods~\cite{chen2018federated, fallah2020personalized, jiang2019improving, khodak2019adaptive,  kulkarni2020survey, nichol2018first} can be extended to customize personalized models by adapting a well-trained global model on the local data of a client~\cite{kairouz2019advances}.
Model mixture methods~\cite{deng2020adaptive, hanzely2020federated} customize for each client by combining the global model with the client's latent local model.
SCAFFOLD~\cite{karimireddy2019scaffold} customizes the gradient updates of personalized models to correct client-drifts between personalized models and a global model.

Most existing local customization methods use a single global model to conduct a global collaboration involving all clients. The global collaboration framework only allows contributions from all clients to a global model and customization of the global model for each client.  It does not allow pairwise collaboration among clients with similar data, and thus may meet dramatic difficulty on non-IID data.

%This universal collaboration can only distinguish the amount of contribution made by each client, but cannot describe the pair-wise collaboration relationships between clients.
%As a result, the local customization methods mentioned above cannot effectively discover the underlying collaboration relationships between different pairs of clients, which largely limits the effectiveness of inter-client collaboration when the data is non-IID across different clients. 

\nop{
However, the customized gradient updates is incompatible with advanced gradient methods, such as ADAM~\cite{kingma2014adam} and AMSGRAD~\cite{reddi2018convergence}. This reduces the performance of SCAFFOLD when clients adopt deep neural networks as personalized models.
}

\nop{
For local customization methods mentioned above, the flexibility of personalized local models are limited because all the local models are customized from the same global model. }

\nop{
which further reduces the final performance of personalized federated learning.
}

\nop{
They can only distinguish the amount of contribution made by each client in the universal collaboration, but cannot discover collaboration groups of clients by modeling the fine-grained pair-wise collaboration relationships between clients.
}

\nop{
can only distinguish the importance of each client in the  collaboration, but cannot model the underlying pairwise 

as the hub to conduct inter-client collaboration can only model the importance of each client in the collaboration, 

but cannot model the underlying pairwise collaboration relationships between clients.

This largely limits the effectiveness of inter-client collaboration, especially when the data is non-IID across different clients, which further impairs the quality of the global model and inevitably reduces the final performance of personalized federated learning.
}



\nop{
makes it difficult to discover the underlying collaboration relationships between each pair of clients, 
}

\nop{
the performance of local customization largely depends on a single well-trained global model.
However, since these methods do not actively discover the underlying collaboration relationships between each pair of clients, the global model is usually trained by taking equal or weighted contributions from all clients, which is difficult to 

since the private data distributions of the clients are unknown, the global model is usually trained by taking equal or weighted contributions from all clients without considering their underlying collaboration relationships.
This inevitably reduces the effectiveness of inter-client collaboration, and further degenerates the performance of personalized models when the data is non-IID.

However, since there is only a single global model that is usually trained by taking equal or weighted contributions from all clients without considering their underlying collaboration relationships, it is difficult to obtain a well-trained global model, especially when the data is non-IID across different clients.
This inevitably reduces the effectiveness of inter-client collaboration, and further degenerates the performance of personalized models.
}

\nop{
More often than not, clients are organized to collaboratively make equal or weighted contributions to train a global model without considering the underlying collaboration relationships between clients.
}

\nop{
the global model is usually trained by aggregating the contributions from all clients with weights related.

without knowing the private data distribution of every client,


letting the clients to contribute 




But it is usually difficult to find a good global model due to the non-IID data of clients~\cite{kairouz2019advances}.


However, these methods\nop{ are effective and easy to implement, but they } can be sensitive to the number of training rounds due to catastrophic forgetting~\cite{mccloskey1989catastrophic, french1999catastrophic} and over-fitting~\cite{deng2020adaptive}.
 


Meta learning methods~\cite{fallah2020personalized, nichol2018first, chen2018federated, khodak2019adaptive, jiang2019improving, kulkarni2020survey} build a global model for multiple clients and produce personalized models by adapting the global model on different clients. These methods mostly focus on synthetic image classification problems, but their performance on the real non-IID data in the scenario of federated learning are yet to be analyzed~\cite{kairouz2019advances}.
}

\nop{
These methods have a good potential to be extended to conduct federated learning~\cite{jiang2019improving, kulkarni2020survey}, however, most existing methods are focusing on synthetic image classification problems, and their performance on the real non-IID data in the scenario of federated learning are yet to be analyzed~\cite{kairouz2019advances}.
}

\nop{
Moreover, requiring the personalized model to be a linear combination of two models largely limits its expressive power.
}

\nop{
can be extended to conduct federated learning by first building a global model for multiple clients and then producing personalized models by adapting the global model on different clients.
}

\nop{
These methods seeks an explicit trade-off between the global model and local models, but 
}

\nop{
let each client to learn a personalized model by a weighted linear combination of the global model and its own local model.
}

\nop{
Model mixture methods explicitly model the trade-off between the global model and local models by letting each client to learn a personalized model by a weighted linear combination of the global model and its own local model.
}




\nop{
\begin{itemize}
    \item Fine tune and domain adaptation~\cite{mansour2009domain, cortes2014domain, ben2010theory} for better personalization
    \item Jiang \textit{et~al.}~\cite{jiang2019improving} highlighted the fact that models of the same structure and performance, but trained differently, can have very different capacity to personalize. In particular, it appears that training models with the goal of maximizing global performance might actually hurt the model’s capacity for subsequent personalization.
    \item The observed gap between the global and personalized accuracy~\cite{jiang2019improving} creates a good argument that personalization should be of central importance to FL.
    \nop{
    \item \chu{All these methods uses a global model for all clients, which inevitably lead to degenerated performance due to the non-IID data distribution on different clients.}
    }
\end{itemize}

\begin{itemize}
    \item Local fine tuning methods involve direct fine tuning methods and meta learning methods. Slightly modify a global model by fine tuning. Two problems, the global model may not be good for non-IID data, and fine tuning may overfit local data. 
    \chu{[Modify] The main drawback of local fine-tuning is that it minimizes the optimization error, whereas the more important part is the generalization performance of the personalized model. In this setting, the personalized model is pruned to overfit~\cite{deng2020adaptive}.}
    \item Model mixing. Two problems, the global model may not be good for non-IID data, and a linear mixture of a global model and a local model limits the expressive power of the mixed model, which makes it difficult to utilize potential clustering structures between clients.
\end{itemize}
}

\nop{
For all these methods, their performance largely relies on the quality of the global model. The need of a global model largely limits their flexibility. Under non-IID data distribution, it may be difficult or even impossible to find a good enough global model to fine tune with.
}

%\textbf{Multi-Task Learning.}
\nop{
Multi-task learning aims to train personalized models for multiple related tasks in a centralized manner by either assuming the structure of the relationships between tasks to be known apriori~\cite{argyriou2007multi, chen2011integrating, evgeniou2004regularized, kim2009statistical}, or attempting to learn the relationships between tasks from data~\cite{gonccalves2016multi, jacob2009clustered, zhang2012convex}.
These methods focus on learning personalized models for different tasks, but ignore the data privacy of clients. Therefore, they are not applicable to federated learning.
To tackle these challenges, Smith \textit{et~al.}~\cite{smith2017federated} extend distributed multi-task learning to federated learning by a primal-dual optimization method,\nop{that is robust to dropped clients and straggling clients}which, however, is not applicable to deep neural networks because strong duality is no longer guaranteed.
}

Smith \textit{et~al.}~\cite{smith2017federated} model the pair-wise collaboration relationships between clients by extending distributed multi-task learning to federated learning. They tackle the problem by a primal-dual optimization method that achieves great performance on convex models. At the same time, due to its rigid requirement of strong duality, their method is not applicable when clients adopt deep neural networks as personalized models.


\nop{Distributed multi-task learning~\cite{vanhaesebrouck2017decentralized, bellet2017personalized, zantedeschi2019fully, xie2017privacy} tried to conduct multi-task learning when the data for each task is separately owned by different clients.
While these methods have the potential to protect the data privacy of clients, they cannot robustly handle the unreliable operating environment of federated learning~\cite{smith2017federated}.
}


\nop{
Vanhaesebrouck \textit{et~al.}~\cite{vanhaesebrouck2017decentralized} and Bellet \textit{et~al.}~\cite{bellet2017personalized} trained personalized models by propagating and averaging models between clients using a pre-exisitng peer-to-peer collaboration network.
However, the collaboration network is usually not known apriori in practice. Zantedeschi~\cite{zantedeschi2019fully} proposed a boosting method to jointly learn the collaboration network with the personalized models of clients, but it needs to pre-train base classifiers on a third-party training data, which is hard to construct without knowing the data distributions of clients.

Xie \textit{et~al.}~\cite{xie2017privacy} proposed a distributed proximal gradient algorithm to train personalized models with guaranteed differential data privacy, but it could not robustly handle client drops.
}

Different from all existing work, our study explores pairwise collaboration among clients.  Our method is particularly effective when clients' data are non-IID, and can take the great advantage of similarity among clients.



\nop{
Most multi-task learning framework tried to model symmetric relationships between multiple tasks (i.e., clients).
Under non-IID setting, the relationship between clients can be asymmetric, which cannot be effectively modeled using classical multi-task learning methods. 
In comparison, attentive message passing does not have any symmetric assumption on the relationship, therefore can perfectly model asymmetric relationships.
}

\nop{Classical multi-task learning can be ported to federated learning.
But they need a $k$, need to inverse a matrix (inefficient).}

\nop{
\begin{itemize}
    \nop{
    \item The general form of multi-task learning is naturally suitable to learn personalized models \chu{cite MOCHA}.
    }
    \nop{
    \item The general form of multi-task learning is also extended to conduct decentralized joint learning of personalized models. These methods encourage similar clients to share more weights and make clients with more data to be more important, but they do not have the attentive mechanism. \chu{These methods need auxiliary information to pre-compute the symmetric relationship between clients, however, such information is not always available, and the collaboration relationships between clients are not always symmetric.}
    }
    \item MOCHA~\cite{smith2017federated}: For example Smith \textit{et~al.} [28] propose a commnication efficient primal-dual optimization method that learns separate but related models for each device through a multi-task learning framework. Despite the theoretical guarantees and practical efficiency of the proposed method, such an approach is not generalizable to non-convex problems, e.g., deep learning, where strong duality is no longer guaranteed.
\end{itemize}
}
% ====================================================
% ====================  Algorithm ====================
% ====================================================





% Consider a federated learning environment with $m$ clients. Suppose that for every $i=1,2,\dots,m$, client $i$ owns $n_i$ samples of data, denoted by $\{z_i^t\}_{t=1}^{n_i}$, that are drawn from a client-specific distribution $\mP_i$. 
% Some notation:
% \begin{itemize}
%     \item $m$: number of clients;
%     \item $n_i$: number of data samples on client $i$, where $i=1,2,\dots,m$;
%     \item $d$: number of parameters in the neural network model;
%     \item $w_i$: model of client $i$, dimension is $d\times 1$; it is stored on client $i$.
%     \item $W$: a matrix of size $d\times m$, whose $i$th column is $w_i$.
%     \item $u_i$: message from client $i$ for collaborative learning; it is stored on the cloud.
%     \item $U$: a matrix of size $d\times m$, whose $i$th column is $u_i$.
% \end{itemize}

% \subsection{An Optimization Formulation}
% We consider the following optimization formulation:
% \begin{equation}
%     \label{eq:form-only-W}
%     \min_{W\in\R^{d\times m}} \mF(W):=\sum_{i=1}^m F_i(w_i) + \lambda\cdot\sum_{\begin{subarray}{c} i,j=1 \\ i\leq j \end{subarray}}^nD(w_i,w_j)
% \end{equation}
% where $\lambda$ is a positive constant, $F_i$ is the empirical loss function of client $i$, i.e.,
% $$ F_i(w) = \frac{1}{n_i}\sum_{j=1}^{n_i} \ell(w; z_j^i), $$
% where $\ell$ is the loss function of the learning task and $\{z_j^i:\, j=1,\dots, n_i\}$ is the data owned by client $i$. Moreover, $D:\R^d\times\R^d\rightarrow \R$ is a non-negative distance function. Examples of $D$ can be:
% \begin{align}
%     D(x,y) & = \frac{1}{2}\|x - y\|^2, \label{eq:ec-dist} \\
%     D(x,y) & = 1 - e^{-\frac{\|x - y\|^2}{2\sigma}}, \label{eq:Gauss-dist}
% \end{align}
% which are the Euclidean distance and Gauss Kernel distance, respectively. To use \eqref{eq:form-only-W} in the setting of federated learning, we propose to apply the following reformulation:
% \begin{equation}
%     \label{eq:form-U=W}
%     \min_{W\in\R^{d\times m}} \mF(W):=\sum_{i=1}^m F_i(w_i) + \lambda\cdot\sum_{\begin{subarray}{c} i,j=1 \\ i\leq j \end{subarray}}^nD(u_i,u_j), \qquad \mbox{s.t.} \ \ U = W,
% \end{equation}
% and employ the quadratic penalty approach to solve \eqref{eq:form-U=W}:
% \begin{equation}
%     \label{eq:form-U-W}
%     \min_{\begin{subarray}{c} U\in\R^{d\times m} \\ W\in\R^{d\times m}\end{subarray}} \mF(W,U) := \sum_{i=1}^m F_i(w_i) + \lambda\cdot\sum_{\begin{subarray}{c} i,j=1 \\ i\leq j \end{subarray}}^nD(u_i,u_j) + \frac{\rho}{2}\|W - U\|_F^2,
% \end{equation}
% where $\rho>0$ is a parameter.

% \subsection{Algorithm}
% We propose an attentive massage passing algorithm for federated multi-task learning (FedAMP), which is based on a gradient descent method for solving our formulation \eqref{eq:form-U-W}. To introduce our FedAMP algorithm, We consider the two distance function in \eqref{eq:ec-dist} and \eqref{eq:Gauss-dist}.
% \subsubsection{Euclidean Distance Function}
% The algorithm FedAMP with Euclidean distance function is given in Algorithm \ref{alg:EDF}.
% \begin{algorithm}[t]
% \DontPrintSemicolon
% \SetKwInput{KwInput}{Input}               
% \SetKwInput{KwOutput}{Output} 
% \KwInput{$\lambda>0$, $\rho>0$, set $u_i^{(0)}=0$ for all $i$}
% \KwOutput{the weight vectors $w_i$ and the message vectors $u_i$ for all $i$}             
% \For{$k=1,2,\dots$}{
% Server sends $u_i^{(k-1)}$ to each client \tcp*[f]{Broadcast}\\
% Let $\mR^{(k)}\subset[m]$ be the set of responded clients in the $k$th round \\
% For each client $i\in\mR^{(k)}$, solve the following training problem: \tcp*[f]{Local Training}
% \begin{equation}
%     \label{eq:update-W-EDF}
%     w_i^{(k)} = \arg\min_w \, F_i(w) + \frac{\rho}{2}\|w - u_i^{(k-1)}\|_2^2.
% \end{equation} 
% \\
% Responded clients send the updated $w_i^{(k)}$ back to Server \tcp*[f]{Collect}\\
% Update the cloud vectors from $u_i^{(k-1)}$ to $u_i^{(k)}$ by: \tcp*[f]{Message Passing}\\
% \begin{equation}
% \label{eq:update-U-EDF}
% u_i^{(k)} = \left\{
% \begin{array}{ll}
% \alpha w_i^{(k)} + (1-\alpha) u_c^{(k-1)}, & \mbox{if} \ i\in\mR^{(k)}, \medskip \\
% u_i^{(k-1)}, & \mbox{if} \ i\notin\mR^{(k)},
% \end{array}
% \right.
% \end{equation}
% where
% $$ \alpha = \frac{\rho}{\rho + \lambda m}, \qquad u_c^{(k-1)} = \frac{\sum_{i=1}^m u_i^{(k-1)}}{m}. $$
% }
% \caption{FedAMP using Euclidean Distance Function (with dropped nodes)}
% \label{alg:EDF}
% \end{algorithm}
% Here we provide some remarks:
% \begin{itemize}
%     \item[(i)] For every $i$, $u_i$ is stored on the cloud and shall be passed explicitly to client $i$ (no privacy). For $w_i$, it is stored on client $i$. In addition, when implementing line 5 of Algorithm \ref{alg:EDF}, we don't need to explicitly send $w_i^{(k)}$. Indeed, according to \eqref{eq:update-U-EDF}, client $i$ can sent back a vector 
%     $$ v_i^{(k)} = \rho(u_i^{(k-1)} - w_i^{(k)}) $$
%     to server, and update $u_i^{(k)}$ by
%     $$ u_i^{(k)} = u_i^{(k-1)} - \frac{v_i^{(k)}}{\rho + \lambda m} - \frac{\lambda m (u_i^{(k-1)} - u_c^{(k-1)})}{\rho + \lambda m}. $$
%     One can verify that the above is equivalent with \eqref{eq:update-U-EDF}. This trick may be useful for improving the ability of protecting privacy in our system.
    
%     \item[(ii)] According to \eqref{eq:update-U-EDF}, for non-responded clients, we don't update its $u_i$ in this round.
% \end{itemize}
    

% \subsubsection{Gauss Kernel Distance Function}
% The algorithm for this part is given in Algorithm \ref{alg:EDF}.
% \begin{algorithm}[t]
% \DontPrintSemicolon
% \SetKwInput{KwInput}{Input}               
% \SetKwInput{KwOutput}{Output} 
% \KwInput{$\lambda>0$, $\rho>0$, $\sigma>0$, set $u_i^{(0)}=0$ for all $i$}
% \KwOutput{the weight vectors $w_i$ and the message vectors $u_i$ for all $i$}             
% \For{$k=1,2,\dots$}{
% Server sends $u_i^{(k-1)}$ to each client \tcp*[f]{Broadcast}\\
% Let $\mR^{(k)}\subset[m]$ be the set of responded clients in the $k$th round \\
% For each client $i\in\mR^{(k)}$, solve the following training problem: \tcp*[f]{Local Training}
% \begin{equation}
%     \label{eq:update-W-GKF}
%     w_i^{(k)} = \arg\min_w \, F_i(w) + \frac{\rho}{2}\|w - u_i^{(k-1)}\|_2^2.
% \end{equation} 
% \\
% Responded clients send the updated $w_i^{(k)}$ back to Server \tcp*[f]{Collect}\\
% Update the cloud vectors from $u_i^{(k-1)}$ to $u_i^{(k)}$ by: \tcp*[f]{Message Passing}\\
% \begin{equation}\label{eq:update-U-GKF}
% u_i^{(k)} = \left\{
% \begin{array}{ll}
% \alpha_i w_i^{(k)} + (1-\alpha_i) u_{i,c}^{(k-1)}, & \mbox{if} \ i\in\mR^{(k)}, \medskip \\
% u_i^{(k-1)}, & \mbox{if} \ i\notin\mR^{(k)},
% \end{array}
% \right.
% \end{equation}
% where for every $i=1,2,\dots,m$, 
% % $$ \alpha_i = \frac{\rho}{\rho + \lambda \beta_i}, \quad \beta_i = \frac{1}{\sigma}\sum_{j=1}^m e^{-\frac{\|u_i^{(k-1)} - u_j^{(k-1)}\|^2}{2\sigma}}, \quad u^{(k-1)}_{i,c} = \frac{\sum_{j=1}^m e^{-\frac{\|u^{(k-1)}_i - u^{(k-1)}_j\|^2}{2\sigma}}u^k_j}{\sum_{j=1}^m e^{-\frac{\|u_i^{(k-1)} - u_j^{(k-1)}\|^2}{2\sigma}}}. $$
% $$ \alpha_i = \frac{\rho}{\rho + \frac{\lambda}{\sigma}\sum_{j=1}^m e^{-\frac{\|u_i^{(k-1)} - u_j^{(k-1)}\|^2}{2\sigma}}}, \qquad u^{(k-1)}_{i,c} = \frac{\sum_{j=1}^m e^{-\frac{\|u^{(k-1)}_i - u^{(k-1)}_j\|^2}{2\sigma}}u^{(k-1)}_j}{\sum_{j=1}^m e^{-\frac{\|u_i^{(k-1)} - u_j^{(k-1)}\|^2}{2\sigma}}}. $$
% }
% \caption{FedAMP using Gauss Kernel Distance Function (with dropped nodes)}
% \label{alg:GKF}
% \end{algorithm}
% Here we provide some remarks:
% \begin{itemize}
%     \item[(i)] The update of $u_i$ in Algorithm \ref{alg:GKF} uses client-specific combination parameters $\alpha_i$'s and client-specific averaged centers $u_{i,c}^{(k-1)}$ (see Eq.~\eqref{eq:update-U-GKF}), whereas those in Algorithm \ref{alg:GKF} are the same for all clients (see Eq.~\eqref{eq:update-U-EDF}).
    
%     \item[(ii)] According to \eqref{eq:update-U-EDF}, for non-responded clients, we don't update its $u_i$ in this round.
% \end{itemize}

% \begin{algorithm}[t]
% \DontPrintSemicolon
% \SetKwInput{KwInput}{Input}               
% \SetKwInput{KwOutput}{Output} 
% \KwInput{$\rho>0$, $\alpha>0$, $w_i^{(0)}$ for all $i$ (pre-train)}
% \KwOutput{the weight vectors $w_i$ for all $i$}             
% \For{$k=1,2,\dots$}{
% Server sends $w_i^{(k-1)}$ to each client \tcp*[f]{Broadcast}\\
% Let $\mR^{(k)}\subset[m]$ be the set of responded clients in the $k$th round \\
% For each client $i\in\mR^{(k)}$, solve the following training problem: \tcp*[f]{Local Training}
% \begin{equation}
%     \label{eq:update-W-GKF}
%     u_i^{(k)} = \arg\min_w \, F_i(w) + \frac{\rho}{2}\|w - w_i^{(k-1)}\|_2^2.
% \end{equation} 
% \\
% For each client $i\notin\mR^{(k)}$, let
% $$ u_i^{(k)} = w_i^{(k-1)}. $$
% \\
% All clients send $u_i^{(k)}$ back to Server \tcp*[f]{Collect}\\
% Update to $w_i^{(k)}$ by: \tcp*[f]{Message Passing}\\
% \begin{equation}\label{eq:update-U-GKF}
% w_i^{(k)} = \left\{
% \begin{array}{ll}
% \alpha u_i^{(k)} + (1-\alpha) u_{i,c}^{(k)}, & \mbox{if} \ i\in\mR^{(k)}, \medskip \\
% u_i^{(k)}, & \mbox{if} \ i\notin\mR^{(k)},
% \end{array}
% \right.
% \end{equation}
% where for every $i=1,2,\dots,m$, 
% % $$ \alpha_i = \frac{\rho}{\rho + \lambda \beta_i}, \quad \beta_i = \frac{1}{\sigma}\sum_{j=1}^m e^{-\frac{\|u_i^{(k-1)} - u_j^{(k-1)}\|^2}{2\sigma}}, \quad u^{(k-1)}_{i,c} = \frac{\sum_{j=1}^m e^{-\frac{\|u^{(k-1)}_i - u^{(k-1)}_j\|^2}{2\sigma}}u^k_j}{\sum_{j=1}^m e^{-\frac{\|u_i^{(k-1)} - u_j^{(k-1)}\|^2}{2\sigma}}}. $$
% $$ u^{(k)}_{i,c} = \frac{\sum_{j=1,j\neq i}^m\frac{u_j^{(k)}}{\|u_i^{(k)} - u_j^{(k)}\|}}{\sum_{j=1,j\neq i}^m\frac{1}{\|u_i^{(k)} - u_j^{(k)}\|}}. $$
% }
% \caption{FedAMP using norm function (with dropped nodes)}
% \label{alg:GKF}
% \end{algorithm}

% \subsection{Framework of FedAMP} 
% Consider a federated learning environment with $m$ clients. Suppose that for every $i=1,2,\dots,m$, client $i$ owns $n_i$ samples of data, denoted by $\{z_{i,t}\}_{t=1}^{n_i}$, that are drawn from a client-specific distribution $\mP_i$. Let $F_i:\R^d\rightarrow\R$ be the training objective function of client $i$. As an example, $F_i$ can be an $\ell_2$-regularized empirical loss function, i.e.,
% $$ F_i(w) = \frac{1}{n_i}\sum_{t=1}^{n_i} \ell_i(w;z_{i,t}) + \frac{\beta_i}{2}\|w\|^2 $$
% for some loss function $\ell_i$ and regularization parameter $\beta_i> 0$ that can be specified by client $i$. In order to improve the learning performance of each client by using the knowledge contained in other clients, we consider the following optimization formulation:
% \begin{equation}
% \label{eq:our-form}
% \min_{w_1,\dots,w_m} \ \left\{\mG(w_1,\dots,w_m):= \sum_{i=1}^m F_i(w_i) + \lambda\cdot\sum_{i<j}^mD(w_i,w_j)\right\},
% \end{equation}
% where $w_i\in\R^d$ is the model to be trained for client $i$, $\lambda>0$ is a regularization parameter, and $D(\cdot,\cdot)$ is a distance-like function that induces collaboration among clients. Concrete examples of the function $D$ that lead to an attentive mechanism in collaboration will be provided in Section \ref{sec:instances}. For now, let us simply assume that $D$ is non-negative, symmetric, and continuously differentiable on $\R^d\times\R^d$.\footnote{We say that $D$ is symmetric if $D(x,y) = D(y,x)$ for any $x,y\in\R^d$.} For ease of reference, we write $\mG = \mF + \lambda\mD$, where
% $$ \mF(w_1,\dots,w_m) = \sum_{i=1}^m F_i(w_i), \quad \mD(w_1,\dots,w_m) = \sum_{i<j}^m D(w_i,w_j). $$
% It worths mentioning that formulations in a similar form of \eqref{eq:our-form} have been adopted in multi-task learning (see, {\it e.g.}, \cite{han2015learning}) but most of the existing algorithms for solving \eqref{eq:our-form} cannot tackle the challenges of data privacy, communication efficiency, stragglers, and dropped clients that are common in the setting of federated learning. Besides, the MOCHA method recently proposed in \cite{smith2017federated} can be applied to \eqref{eq:our-form} but is limited to the scenario where every $F_i$ is convex.

% Based on \eqref{eq:our-form}, we propose a novel algorithmic framework named federated attentive message passing ({FedAMP}) that promotes effective collaboration in training and also addresses the aforementioned challenges. It stems from an incremental type of optimization algorithm (see \cite{bertsekas2011incremental}) applied to \eqref{eq:our-form} and it is carefully deployed to meet the demand of federated learning environment. Specifically, {FedAMP} starts with an initial guess of the models $(w_1^{(0)}, \dots, w_m^{(0)})$ that are stored on the server. At the $k$-th communication round of {FedAMP}, the server first computes a set of intermediate models $(u_1^{(k)},\dots,u_m^{(k)})$ by applying a gradient descent step on the regularization function $\mD$ in \eqref{eq:our-form}, {\it i.e.},
% \begin{equation}\label{eq:grad-step}
% u_i^{(k)} = w_i^{(k-1)} - \alpha_{k,i}\nabla_i\mD(w_1^{(k-1)},\dots,w_m^{(k-1)}), \quad \forall i=1,\dots,m,
% \end{equation}
% where $\alpha_{k,i}>0$ is the step size for the $i$-th model and $\nabla_i\mD$ stands for the partial gradient of $\mD$ with respect to $w_i$. Next, the server sends each intermediate model $u_i^{(k)}$ to client $i$ and requests it to perform local training with the knowledge of $u_i^{(k)}$. In particular, client $i$ computes an updated model $w_i^{(k)}$ by
% \begin{equation}
% \label{eq:local-training}
% w_i^{(k)} \approx \arg\min_{w\in\R^d} F_i(w) + \frac{1}{2\alpha_{k,i}}\|w - u_i^{(k)}\|^2,
% \end{equation}
% where $\alpha_{k,i}$ is given in \eqref{eq:grad-step}. Finally, the server collects the updated models $(w_1^{(k)},\dots,w_m^{(k)})$ and the system enters the $(k+1)$-th communication round. When {FedAMP} terminates at the $K$-th round for some $K$, each client $i$ naturally holds a personalized model $w_i^{(K)}$, which is considered to be the output of {FedAMP}. The whole procedure of {FedAMP} is summarized in Algorithm \ref{alg:FedAMP}.

% \begin{algorithm}[t]
% \DontPrintSemicolon
% \SetKwInput{KwInput}{Input}               
% \SetKwInput{KwOutput}{Output} 
% \KwInput{number of clients $m$, initial models $(w_1^{(0)},\dots, w_m^{(0)})$, total communication rounds $K$}
% %\KwOutput{the models $w_i^{(K)}$}             
% \For{$k=1,2,\dots, K$}{
% server computes the intermediate models $(u_1^{(k)}, \dots, u_m^{(k)})$ by \eqref{eq:grad-step} \\
% %\tcp*[f]{Server Computation}
% %\begin{equation}\label{eq:grad-step-c}
% %u_i^{(k)} = w_i^{(k-1)} - \alpha_{k,i}\nabla_i\mD(w_1^{(k-1)},\dots,w_m^{(k-1)}), \quad i=1,\dots,m
% %\end{equation}\\
% server sends each $u_i^{(k)}$ to client $i$ \\
% clients perform local training and compute the updated models $(w_1^{(k)},\dots,w_m^{(k)})$ by \eqref{eq:local-training}\\
% %\begin{equation}
% %\label{eq:local-training-c}
% %w_i^{(k)} \approx \arg\min_{w\in\R^d} F_i(w) + \frac{1}{2\alpha_{k,i}}\|w - u_i^{(k)}\|^2
% %\end{equation}
% server collects the updated models $(w_1^{(k)},\dots,w_m^{(k)})$
% }
% \caption{{FedAMP}: Federated Learning with Attentive Message Passing}
% \label{alg:FedAMP}
% \end{algorithm}

% We make the following remarks on {FedAMP}:
% \begin{itemize}
% \item[(i)] The initial models $(w_1^{(0)}, \dots, w_m^{(0)})$ can be obtained, for example, by separately pre-training on each client with their local data and then collected by the server.

% \item[(ii)] As we shall elaborate in Section \ref{sec:instances}, for many distance-like functions $D$ in \eqref{eq:our-form}, the gradient descent step in \eqref{eq:grad-step} amounts to attentively aggregating the models. In particular, for every $i$, the intermediate model $u_i^{(k)}$ can be written as a convex combination of $(w_1^{(k-1)},\dots,w_m^{(k-1)})$, {\it i.e.}, 
% %$u_i^{(k)} = \xi_{1}w_1^{(k-1)} + \dots + \xi_{m}w_m^{(k-1)}$
% \begin{equation}
% \label{eq:weight-average}
% u_i^{(k)} = \xi_{1}w_1^{(k-1)} + \dots + \xi_{m}w_m^{(k-1)}
% \end{equation}
% for some non-negative weights $\{\xi_{j}\}_{j=1}^m$ such that $\xi_1 + \dots + \xi_m = 1$. Moreover, the magnitude of $\xi_j$ is proportional to the similarity between $w_i^{(k-1)}$ and $w_j^{(k-1)}$, leading to an attentive mechanism in aggregation \eqref{eq:weight-average}. Furthermore, the mechanism in aggregation is blind to the clients, as the distance-like function $D$ is chosen by the server and does not need to be shared with the clients.

% \item[(iii)] One can observe that we require an exact computation in \eqref{eq:grad-step} but allow an approximate solution in \eqref{eq:local-training}. This is due to the fact that the computation in \eqref{eq:grad-step} is performed on the server, which is stable and can be well controlled, while the computation in \eqref{eq:local-training} is performed across the heterogeneous networks, where stragglers or even dropped clients are common. Details on the approximation criterion for \eqref{eq:local-training} shall be specified in Section \ref{sec:analysis} when we analyze the convergence of the FedAMP framework. In practice, one may simply apply a number of epochs of stochastic optimization algorithms to obtain an updated model $w_i^{(k)}$ in \eqref{eq:local-training}. 
% \end{itemize} 




% \begin{align}
%     D(x,y) & = \frac{1}{2}\|x - y\|^2, \label{eq:ec-dist} \\
%     D(x,y) & = \|x - y\|, \label{eq:norm-dist} \\
%     D(x,y) & = 1 - e^{-\frac{\|x - y\|^2}{2\sigma}}, \label{eq:Gauss-dist}
% \end{align}


% \subsection{Examples of Attentive Collaboration}\label{sec:instances}
% Below we consider three examples of the distance-like function $D$ 
% We propose an attentive massage passing algorithm for federated multi-task learning (FedAMP), which is based on a gradient descent method for solving our formulation \eqref{eq:form-U-W}. To introduce our FedAMP algorithm, We consider the two distance function in \eqref{eq:ec-dist} and \eqref{eq:Gauss-dist}.

% \subsection{Convergence Analysis}\label{sec:analysis}
    
% \subsection{Convergence Guarantee}
% We may also consider the following formulation:
% \begin{equation}
%     \label{eq:form-only-U}
%     \min_{U\in\R^{d\times m}} \ \mH(U): = \underbrace{\sum_{i=1}^m g_i(u_i)}_{\mG(U)} + \lambda \cdot \underbrace{\sum_{\begin{subarray}{c} i,j=1 \\ i\leq j \end{subarray}}^n D(u_i,u_j)}_{\mD(U)},
% \end{equation}
% where $g_i:\R^d\rightarrow\R$ is defined as
% \begin{equation}
%     \label{eq:-def-Gi}
%     g_i(u) := \min_{w\in\R^d} F_i(w) + \frac{\rho}{2}\|w - u\|_2^2
% \end{equation}
% with $F_i:\R^d\rightarrow\R$ being the empirical loss function of Client $i$. Then, we can apply gradient descent method to solve \eqref{eq:form-only-U}:
% \begin{equation}
%     \label{eq:grad-descent}
%     U^{k+1} = U^k - \alpha_k\nabla \mH(U^k), \quad k=0,1,2,\dots.
% \end{equation}
% Now we show how to compute $\nabla \mH(U^k)$ under the condition that every $F_i$ is convex. First, for every $g_i$, we have
% $$\nabla g_i(u^k_i) = \rho(u^k_i - w^k_i), $$
% where
% $$ w^k_i = \arg\min_{w\in\R^d} F_i(w) + \frac{\rho}{2}\|w - u^k_i \|_2^2. $$
% Thus, we obtain
% \begin{equation}
%     \label{eq:nabla-G}
%     \nabla \mG(U^k) = \rho\cdot\left[u^k_1 - w^k_1, \ u^k_2 - w^k_2, \ \dots, \ u^k_m - w^k_m\right].
% \end{equation}
% Moreover, by theory of convex analysis, we know that $\nabla \mG$ is Lipschitz continuous with Lipschitz constant $\rho$.

% Next, we consider the two distance functions in \eqref{eq:ec-dist} and \eqref{eq:Gauss-dist} separately.

% \subsubsection{Euclidean Distance Function}
% In the case of \eqref{eq:ec-dist}, one can verify that 
% \begin{equation}\label{eq:nabla-D-1}
% % \nabla\left(\sum_{\begin{subarray}{c} i,j=1 \\ i\leq j \end{subarray}}^n D(u^k_i,u^k_j)\right)
% \nabla\mD(U^k)
% = m\left[ u^k_1 - u^k_c, \ u^k_2 - u^k_c, \ \dots, \ u^k_m - u^k_c \right], \quad \mbox{where} \quad u^k_c = \frac{\sum_{i=1}^mu^k_i}{m}.
% \end{equation}
% Equivalently, we can write
% $$ \nabla \mD(U^k) = mU^k - U^k\mathbf{E}, $$
% where $\mathbf{E}$ is the all-one $m\times m$ matrix. Thus, for any $U_1$ and $U_2$, we have
% $$ \|\nabla \mD(U_1) - \nabla \mD(U_2\|_F \leq \|U_1 - U_2\|_F\|mI - \mathbf{E}\|_2 = m\|U_1 - U_2\|_F, $$
% which implies that $\nabla \mD$ is Lipschitz continuous with Lipschitz constant $m$. This, together with \eqref{eq:nabla-G} and \eqref{eq:form-only-U}, yields that $\nabla\mH$ has Lipschitz constant $\rho + \lambda m$. Hence, we can take $\alpha_k\equiv(\rho+\lambda m)^{-1}$ in \eqref{eq:grad-descent}.

% By \eqref{eq:grad-descent}, \eqref{eq:nabla-G}, and \eqref{eq:nabla-D-1}, and taking $\alpha_k \equiv 1/(\rho+\lambda m)$, we can write the update of $U$ as
% $$ 
% \begin{aligned}
% u^{k+1}_i & = u^k_i - \alpha_k\big(\rho(u^k_i - w^k_i) + \lambda m(u^k_i - u^k_c)\big) = (1 - \alpha_k(\rho+\lambda m))u_i^k + \alpha_k\rho w_i^k + \alpha_k \lambda m u_c^k \\ 
% & = \frac{\rho}{\rho+\lambda m}w_i^k + \frac{\lambda m}{\rho+\lambda m}u_c^k.
% \end{aligned}
% $$
% Notice that this is almost identical to our heuristic method.
% \subsubsection{Gaussian Kernel Distance Function}
% In the case of \eqref{eq:Gauss-dist}, one can verify that
% \begin{equation}\label{eq:nabla-D-2}
% % \nabla\left(\sum_{\begin{subarray}{c} i,j=1 \\ i\leq j \end{subarray}}^n D(u^k_i,u^k_j)\right)=
% \nabla\mD(U^k) = \left[ c_1^k(u^k_1 - u^k_{1,c}), \ c_2^k(u^k_2 - u^k_{2,c}), \ \dots, \ c_m^k(u^k_m - u^k_{m,c})\right],
% \end{equation}
% where
% $$ c_i^k = \frac{1}{\sigma}\sum_{j=1}^m e^{-\frac{\|u^k_i - u^k_j\|^2}{2\sigma}}, \quad u^k_{i,c} = \frac{\sum_{j=1}^m e^{-\frac{\|u^k_i - u^k_j\|^2}{2\sigma}}u^k_j}{\sigma c_i^k}, \quad i=1,2,\dots,m.$$
% Some preliminary ideas show that $\nabla\mD$ may also be Lipschitz continuous with constant $m$. Then, similar to the previous case, we can take $\alpha_k\leq (\rho+\lambda m)^{-1}$.

% By \eqref{eq:grad-descent}, \eqref{eq:nabla-G}, and \eqref{eq:nabla-D-2}, we can write the update of $U$ as
% $$ 
% \begin{aligned}
% u^{k+1}_i & = u_i^k - \alpha_k\big(\rho(u_i^k - w_i^k) + \lambda c_i^k(u_i^k - u_{i,c}^k)\big) \\
% & = (1 - \alpha_k\rho - \alpha_k\lambda c_i^k)u_i^k + \alpha_k\rho w_i^k + \alpha_k\lambda c_i^k u_{i,c}^k \\
% & = (1 - \alpha_k\rho - \alpha_k\lambda c_i^k)u_i^k + \alpha_k\rho w_i^k + \alpha_k\lambda c_i^k\cdot \frac{\sum_{j=1}^m e^{-\frac{\|u^k_i - u^k_j\|^2}{2\sigma}}u^k_j}{\sigma c_i^k}
% \end{aligned}
% $$
% If we choose $\alpha_k = (\rho + \lambda c_i^k)^{-1}$, then
% \begin{equation}
%     \label{eq:U-update-gauss}
%     \begin{aligned}
%     u_i^{k+1} & = \frac{\rho}{\rho + \lambda c_i^{k}}w_i^k + \frac{\lambda c_i^k}{\rho + \lambda c_i^k}\frac{\sum_{j=1}^m e^{-\frac{\|u^k_i - u^k_j\|^2}{2\sigma}}u^k_j}{\sum_{j=1}^m e^{-\frac{\|u^k_i - u^k_j\|^2}{2\sigma}}} \\
%     & = \frac{\rho}{\rho + \lambda c_i^{k}}w_i^k + \frac{\lambda }{\sigma(\rho + \lambda c_i^k)}\sum_{j=1}^m e^{-\frac{\|u^k_i - u^k_j\|^2}{2\sigma}}u^k_j 
%     \end{aligned}
% \end{equation} 
% Almost identical to our implementation.

% {\bf New:} Just do gradient descent for $U$, with step size $\alpha>0$:
% \begin{align*}
% w_i^{(k+1)} & = u_i^{(k)} - \alpha\nabla_i D(U^{(k)}) \\
% & = u_i^{(k)} - \alpha\cdot \left(c_i^{(k)}(u_i^{(k)} - u_{i,c}^{(k)})\right) \\
% & = u_i^{(k)} - \frac{\alpha}{\sigma}\cdot\sum_{j=1}^m e^{-\frac{\|u_i^{(k)} - u_j^{(k)}\|^2}{2\sigma}}\left(u_i^{(k)} - u_j^{(k)}\right) \\
% & = u_i^{(k)} - \frac{\alpha}{\sigma}\cdot\sum_{j=1,j\neq i}^m e^{-\frac{\|u_i^{(k)} - u_j^{(k)}\|^2}{2\sigma}}\left(u_i^{(k)} - u_j^{(k)}\right) \\
% & = \left(1 - \frac{\alpha}{\sigma}\cdot\sum_{j=1,j\neq i}^m e^{-\frac{\|u_i^{(k)} - u_j^{(k)}\|^2}{2\sigma}}\right)u_i^{(k)} + \frac{\alpha}{\sigma}\cdot\sum_{j=1,j\neq i}^m e^{-\frac{\|u_i^{(k)} - u_j^{(k)}\|^2}{2\sigma}}u_j^{(k)}.
% \end{align*}
% Let 
% $$ \alpha = \beta\sigma\cdot\left(\sum_{j=1,j\neq i}^m e^{-\frac{\|u_i^{(k)} - u_j^{(k)}\|^2}{2\sigma}}\right)^{-1} $$
% for some $\beta>0$. Then,
% $$ w_i^{(k+1)} = (1-\beta)u_i^{(k)} + \beta\cdot \frac{\sum_{j=1,j\neq i}^m e^{-\frac{\|u_i^{(k)} - u_j^{(k)}\|^2}{2\sigma}}u_j^{(k)}}{\sum_{j=1,j\neq i}^m e^{-\frac{\|u_i^{(k)} - u_j^{(k)}\|^2}{2\sigma}}}.$$

% \subsection{Framework of FedAMP} 
% Consider a federated learning environment with $m$ clients. Suppose that for every $i=1,2,\dots,m$, client $i$ owns $n_i$ samples of data, denoted by $\{z_{i,t}\}_{t=1}^{n_i}$, that are drawn from a client-specific distribution $\mP_i$. Let $F_i:\R^d\rightarrow\R$ be the training objective function of client $i$. As an example, $F_i$ can be an $\ell_2$-regularized empirical loss function, i.e.,
% $$ F_i(w) = \frac{1}{n_i}\sum_{t=1}^{n_i} \ell_i(w;z_{i,t}) + \frac{\gamma_i}{2}\|w\|^2 $$
% for some loss function $\ell_i$ and regularization parameter $\gamma_i> 0$ that can be specified by client $i$. In order to improve the learning performance of each client by using the knowledge contained in other clients, we consider the following optimization formulation:
% \begin{equation}
% \label{eq:our-form}
% \min_{w_1,\dots,w_m} \ \left\{\mG(w_1,\dots,w_m):= \sum_{i=1}^m F_i(w_i) + \lambda\cdot\sum_{i<j}^mD(w_i,w_j)\right\},
% \end{equation}
% where $w_i\in\R^d$ is the model to be trained for client $i$, $\lambda>0$ is a regularization parameter, and $D(\cdot,\cdot)$ is a distance-like function that induces collaboration among clients. Concrete examples of the function $D$ that lead to an attentive mechanism in collaboration will be provided in Section \ref{sec:instances}. For now, let us simply assume that $D$ is non-negative, symmetric, and continuously differentiable on $\R^d\times\R^d$.\footnote{We say that $D$ is symmetric if $D(x,y) = D(y,x)$ for any $x,y\in\R^d$.} For ease of reference, we write $\mG = \mF + \lambda\mD$, where
% $$ \mF(w_1,\dots,w_m) = \sum_{i=1}^m F_i(w_i), \quad \mD(w_1,\dots,w_m) = \sum_{i<j}^m D(w_i,w_j). $$
% It worths mentioning that formulations in a similar form of \eqref{eq:our-form} have been adopted in multi-task learning (see, {\it e.g.}, \cite{han2015learning}) but most of the existing algorithms for solving \eqref{eq:our-form} cannot tackle the challenges of data privacy, communication efficiency, stragglers, and dropped clients that are common in the setting of federated learning. Besides, the {\sc Mocha} method recently proposed in \cite{smith2017federated} can be applied to \eqref{eq:our-form} but is limited to the scenario where every $F_i$ is convex.

% Based on \eqref{eq:our-form}, we propose a novel algorithmic framework named federated attentive message passing ({FedAMP}) that promotes effective collaboration in training and also addresses the aforementioned challenges. It stems from an incremental type of optimization algorithm (see \cite{bertsekas2011incremental}) applied to \eqref{eq:our-form} and it is carefully deployed to meet the demand of federated learning environment. Specifically, {FedAMP} starts with an initial guess of the models $(w_1^{(0)}, \dots, w_m^{(0)})$ that are stored on the server. At the $k$-th communication round of {FedAMP}, the server first computes a set of intermediate models $(u_1^{(k)},\dots,u_m^{(k)})$ by applying a gradient descent step on the regularization function $\mD$ in \eqref{eq:our-form}, {\it i.e.},
% \begin{equation}\label{eq:grad-step}
% u_i^{(k)} = w_i^{(k-1)} - \alpha_{k,i}\nabla_i\mD(w_1^{(k-1)},\dots,w_m^{(k-1)}), \quad \forall i=1,\dots,m,
% \end{equation}
% where $\alpha_{k,i}>0$ is the step size for the $i$-th model and $\nabla_i\mD$ stands for the partial gradient of $\mD$ with respect to $w_i$. Next, the server sends each intermediate model $u_i^{(k)}$ to client $i$ and requests it to perform local training with the knowledge of $u_i^{(k)}$. In particular, client $i$ computes an updated model $w_i^{(k)}$ by
% \begin{equation}
% \label{eq:local-training}
% w_i^{(k)} \approx \arg\min_{w\in\R^d} F_i(w) + \frac{1}{2\beta_{k,i}}\|w - u_i^{(k)}\|^2
% \end{equation}
% for some $\beta_{k,i}>0$. Finally, the server collects the updated models $(w_1^{(k)},\dots,w_m^{(k)})$ and the system enters the $(k+1)$-th communication round. When {FedAMP} terminates at the $K$-th round for some $K$, each client $i$ naturally holds an intermediate model $u_i^{(k)}$ and a personalized model $w_i^{(K)}$, which are considered to be the output of {FedAMP}. The whole procedure of {FedAMP} is summarized in Algorithm \ref{alg:FedAMP}.

% \begin{algorithm}[t]
% \DontPrintSemicolon
% \SetKwInput{KwInput}{Input}               
% \SetKwInput{KwOutput}{Output} 
% \KwInput{number of clients $m$, initial models $(w_1^{(0)},\dots, w_m^{(0)})$, total communication rounds $K$, gradient descent step sizes $\{\alpha_{k,i}\}\subset\R_{++}$, local training parameters $\{\beta_{k,i}\}\subset\R_{++}$}
% \KwOutput{the final models $(w_1^{(K)},\dots, w_m^{(K)})$ and $(u_1^{(K)},\dots, u_m^{(K)})$}             
% \For{$k=1,2,\dots, K$}{
% server computes the intermediate models $(u_1^{(k)}, \dots, u_m^{(k)})$ by \eqref{eq:grad-step} \\
% %\tcp*[f]{Server Computation}
% %\begin{equation}\label{eq:grad-step-c}
% %u_i^{(k)} = w_i^{(k-1)} - \alpha_{k,i}\nabla_i\mD(w_1^{(k-1)},\dots,w_m^{(k-1)}), \quad i=1,\dots,m
% %\end{equation}\\
% server sends each $u_i^{(k)}$ to client $i$ \\
% clients perform local training and compute the updated models $(w_1^{(k)},\dots,w_m^{(k)})$ by \eqref{eq:local-training}\\
% %\begin{equation}
% %\label{eq:local-training-c}
% %w_i^{(k)} \approx \arg\min_{w\in\R^d} F_i(w) + \frac{1}{2\alpha_{k,i}}\|w - u_i^{(k)}\|^2
% %\end{equation}
% server collects the updated models $(w_1^{(k)},\dots,w_m^{(k)})$
% }
% \caption{Federated Learning with Attentive Message Passing ({FedAMP})}
% \label{alg:FedAMP}
% \end{algorithm}

% We make the following remarks on {FedAMP}:
% \begin{itemize}
% \item[(i)] The initial models $(w_1^{(0)}, \dots, w_m^{(0)})$ can be obtained, for example, by separately pre-training on each client with their local data and then collected by the server.

% \item[(ii)] As we shall elaborate in Section \ref{sec:instances}, for many distance-like functions $D$ in \eqref{eq:our-form}, the gradient descent step in \eqref{eq:grad-step} amounts to attentively aggregating the models. In particular, for every $i$, the intermediate model $u_i^{(k)}$ can be written as a convex combination of $(w_1^{(k-1)},\dots,w_m^{(k-1)})$, {\it i.e.}, 
% %$u_i^{(k)} = \xi_{1}w_1^{(k-1)} + \dots + \xi_{m}w_m^{(k-1)}$
% \begin{equation}
% \label{eq:weight-average}
% u_i^{(k)} = \xi_{1}w_1^{(k-1)} + \dots + \xi_{m}w_m^{(k-1)}
% \end{equation}
% for some non-negative weights $\{\xi_{j}\}_{j=1}^m$ such that $\xi_1 + \dots + \xi_m = 1$. Moreover, the magnitude of $\xi_j$ is proportional to the similarity between $w_i^{(k-1)}$ and $w_j^{(k-1)}$, leading to an attentive mechanism in aggregation \eqref{eq:weight-average}. Furthermore, the mechanism in aggregation is blind to the clients, as the distance-like function $D$ is chosen by the server and does not need to be shared with the clients.

% \item[(iii)] One can observe that we require an exact computation in \eqref{eq:grad-step} but allow an approximate solution in \eqref{eq:local-training}. This is due to the fact that the computation in \eqref{eq:grad-step} is performed on the server, which is stable and can be well controlled, while the computation in \eqref{eq:local-training} is performed across the heterogeneous networks, where stragglers or even dropped clients are common. Details on the approximation criterion for \eqref{eq:local-training} shall be specified in Section \ref{sec:analysis} when we analyze the convergence of the FedAMP framework. In practice, one may simply apply a number of epochs of stochastic optimization algorithms to obtain an updated model $w_i^{(k)}$ in \eqref{eq:local-training}. 
% \end{itemize} 

% %
% %
% %
% %\begin{align}
% %    D(x,y) & = \frac{1}{2}\|x - y\|^2, \label{eq:ec-dist} \\
% %    D(x,y) & = \|x - y\|, \label{eq:norm-dist} \\
% %    D(x,y) & = 1 - e^{-\frac{\|x - y\|^2}{2\sigma}}, \label{eq:Gauss-dist}
% %\end{align}


% \subsection{Attentive Mechanisms in Collaboration}\label{sec:instances}
% In this subsection, we elaborate on the collaboration among clients in {FedAMP}. In particular, we instantiate the {FedAMP} framework by three concrete examples of the distance-like function $D$, with which the gradient descent step in \eqref{eq:grad-step} promotes an attentive mechanism in collaboration.

% \subsubsection{Squared Euclidean Distance}
% As a warm-up, we first instantiate the {FedAMP} framework with $D$ being the squared Euclidean distance:
% \begin{equation}
% \label{eq:s-ed}
% D(x,y) = \frac{1}{2}\|x - y\|^2,
% \end{equation}
% where $\|\cdot\|$ stands for the $\ell_2$-norm in $\R^d$. Upon substituting \eqref{eq:s-ed} into \eqref{eq:our-form}, one can easily verify that
% $$ \nabla_i\mD(w_1,\dots,w_m) = \sum_{j\neq i}^m (w_i - w_j). $$
% It then follows from \eqref{eq:grad-step} that
% $$ u_i^{(k)} = w_i^{(k-1)} - \alpha_{k,i}\sum_{j\neq i}^m(w_i^{(k-1)} - w_j^{(k-1)}) = (1 - (m-1)\alpha_{k,i})w_i^{(k-1)} + \alpha_{k,i}\sum_{j\neq i}^mw_j^{(k-1)}. $$
% In light of this, we choose $\alpha_{k,i} = \tau_{k,i}/(m-1)$ for some $\tau_{k,i}\in(0,1]$, which leads to the following simple implementation of \eqref{eq:grad-step}:
% \begin{equation}
% \label{eq:gs-s-ed}
% u_i^{(k)} = (1 - \tau_{k,i})\cdot w_i^{(k-1)} + \tau_{k,i}\cdot\frac{\sum_{j\neq i}^mw_j^{(k-1)}}{m-1}.
% \end{equation}
% This yields that $u_i^{(k)}$ is a weighted average of $(w_1^{(k-1)},\dots,w_m^{(k-1)})$ with a parameter $\tau_{k,i}$ controlling the attention to the $i$-th client itself. Moreover, for every $j\neq i$, $w_j^{(k-1)}$ contributes equally to the aggregation \eqref{eq:gs-s-ed}. 

% From \eqref{eq:gs-s-ed} and Algorithm \ref{alg:FedAMP}, we observe that {FedAMP} with squared Euclidean distance function \eqref{eq:s-ed} can be considered as an extension of FedAvg in \cite{li2018federated}. Indeed, by letting $\tau_{k,i}\equiv (m-1)/m$ in \eqref{eq:gs-s-ed}, {FedAMP} reduces to the general form of FedAvg. Moreover, with the parameter $\tau_{k,i}$ in \eqref{eq:gs-s-ed}, {FedAMP} is able to construct personalized intermediate models $(u_1^{(k)},\dots,u_m^{(k)})$ that shall be used as the prox-centers in local training \eqref{eq:local-training}. 

% \subsubsection{Smoothed Euclidean Distance}
% The second instance we would like to consider is the Euclidean distance function $\|x - y\|$. However, since $\|\cdot\|$ is not differentiable at the origin, the Euclidean distance function does not satisfy our assumption for $D$. Instead, we consider a smoothed version of it defined as:
% \begin{equation}
% \label{eq:sm-ed}
% D(x,y) = \left\{
% \begin{aligned}
% & \frac{1}{2\mu}\|x - y\|^2, \quad \ \mbox{if} \ \|x - y\|\leq \mu, \\
% & \|x - y\| - \frac{\mu}{2}, \quad \mbox{otherwise},
% \end{aligned}\right.
% \end{equation}
% where $\mu>0$ is a smoothing parameter. Upon substituting \eqref{eq:sm-ed} into \eqref{eq:our-form}, one can verify that
% %$$ \nabla_i\mD(w_1,\dots,w_m) = \sum_{\begin{subarray}{c}
% %j\neq i \\ \|w_i - w_j\|\leq \mu
% %\end{subarray}} \frac{w_i - w_j}{\mu} + \sum_{\begin{subarray}{c}
% %j\neq i \\ \|w_i - w_j\|> \mu
% %\end{subarray}}\frac{w_i - w_j}{\|w_i - w_j\|}. $$
% $$ \nabla_i\mD(w_1,\dots,w_m) = \sum_{j\neq i}^m \frac{w_i - w_j}{\max\left\{\mu, \, \|w_i - w_j\|\right\}} = \sum_{j\neq i}^m\min\left\{\mu^{-1},\,\|w_i - w_j\|^{-1}\right\}\cdot(w_i - w_j) $$
% It then follows from \eqref{eq:grad-step} that
% $$ u_i^{(k)} = w_i^{(k-1)} - \alpha_{k,i}\sum_{j\neq i}^m \min\left\{\mu^{-1}, \, \|w_i^{(k-1)} - w_j^{(k-1)}\|^{-1}\right\}\cdot\left(w_i^{(k-1)} - w_j^{(k-1)}\right). $$
% Then, by choosing $\alpha_{k,i} = \tau_{k,i}\cdot\left(\sum_{j\neq i}^m \min\left\{\mu^{-1}, \, \|w_i^{(k-1)} - w_j^{(k-1)}\|^{-1}\right\}\right)^{-1}$ for some $\tau_{k,i}\in(0,1]$, the step \eqref{eq:grad-step} can be implemented as:
% \begin{equation}
% \label{eq:gs-sm-ed}
% u_i^{(k)} = (1 - \tau_{k,i})\cdot w_i^{(k-1)} + \tau_{k,i}\cdot\frac{\sum_{j\neq i}^m \min\left\{\mu^{-1}, \, \|w_i^{(k-1)} - w_j^{(k-1)}\|^{-1}\right\}\cdot w_j^{(k-1)}}{\sum_{j\neq i}^m \min\left\{\mu^{-1}, \, \|w_i^{(k-1)} - w_j^{(k-1)}\|^{-1}\right\}}.
% \end{equation}
% Similar to \eqref{eq:gs-s-ed}, here $u_i^{(k)}$ is also a weighted average of $(w_1^{(k-1)},\dots,w_m^{(k-1)})$ with a parameter $\tau_{k,i}$ controlling the attention to the $i$-th client itself. On the other hand, unlike \eqref{eq:gs-s-ed}, the weights of $w_j^{(k-1)}$ for all $j\neq i$ are not equally distributed in \eqref{eq:gs-sm-ed}. Indeed, for any $j\neq i$, the weight of $w_j^{(k-1)}$ in \eqref{eq:gs-sm-ed} is proportional to $\min\left\{\mu^{-1}, \, \|w_i^{(k-1)} - w_j^{(k-1)}\|^{-1}\right\}$. This implies that an $w_j^{(k-1)}$ with $j\neq i$ shall contribute more to the aggregation \eqref{eq:gs-sm-ed} if $\|w_i^{(k-1)} - w_j^{(k-1)}\|$ is smaller, or equivalently, if $w_j^{(k-1)}$ is closer to $w_i^{(k-1)}$.  

% \subsubsection{Gaussian Kernel Distance}
% The third instance we consider is the distance function induced by the Gaussian kernel, which is defined as
% \begin{equation}
% \label{eq:gk-dist}
% D(x,y) = 1 - e^{-\frac{\|x - y\|^2}{2\sigma}}
% \end{equation}
% for some $\sigma>0$. Upon substituting \eqref{eq:gk-dist} into \eqref{eq:our-form}, one can verify that
% $$ \nabla_i\mD(w_1,\dots,w_m) = \frac{1}{\sigma}\sum_{j\neq i}^me^{-\frac{\|w_i - w_j\|^2}{2\sigma}}(w_i - w_j). $$
% It then follows from \eqref{eq:grad-step} that
% $$ u_i^{(k)} = w_i^{(k-1)} - \frac{\alpha_{k,i}}{\sigma}\sum_{j\neq i}^me^{-\frac{\|w_i^{(k-1)} - w_j^{(k-1)}\|^2}{2\sigma}}(w_i^{(k-1)} - w_j^{(k-1)}). $$
% Then, by choosing $\alpha_{k,i} = \tau_{k,i}\sigma\left(\sum_{j\neq i}^me^{-\frac{\|w_i^{(k-1)} - w_j^{(k-1)}\|^2}{2\sigma}}\right)^{-1}$ for some $\tau_{k,i}\in(0,1]$, the step \eqref{eq:grad-step} can be implemented as:
% \begin{equation}
% \label{eq:gs-gk-dist}
% u_i^{(k)} = (1 - \tau_{k,i})\cdot w_i^{(k-1)} + \tau_{k,i}\cdot\frac{\sum_{j\neq i}^m e^{-\frac{\|w_i^{(k-1)} - w_j^{(k-1)}\|^2}{2\sigma}}\cdot w_j^{(k-1)}}{\sum_{j\neq i}^m e^{-\frac{\|w_i^{(k-1)} - w_j^{(k-1)}\|^2}{2\sigma}}}.
% \end{equation}
% Again $u_i^{(k)}$ is a weighted average of $(w_1^{(k-1)},\dots,w_m^{(k-1)})$ with a parameter $\tau_{k,i}$ controlling the attention to the $i$-th client itself. Moreover, for any $j\neq i$, the weight of $w_j^{(k-1)}$ in \eqref{eq:gs-gk-dist} is proportional to $e^{-\frac{\|w_i^{(k-1)} - w_j^{(k-1)}\|^2}{2\sigma}}$. Similar to \eqref{eq:gs-sm-ed}, this implies that an $w_j^{(k-1)}$ with $j\neq i$ shall contribute more to the aggregation \eqref{eq:gs-gk-dist} if $w_j^{(k-1)}$ is closer to $w_i^{(k-1)}$.  

% In view of \eqref{eq:gs-s-ed}, \eqref{eq:gs-sm-ed}, and \eqref{eq:gs-gk-dist}, it is clear that {FedAMP} can promote attentive collaboration among the personalized models $(w_1,\dots,w_m)$. Indeed, the intermediate model $u_i^{(k)}$, which shall be used in the $k$-th round of {FedAMP} as the prox-center to enhance the local training on client $i$, is a convex combination of the models $(w_1^{(k-1)},\dots,w_m^{(k-1)})$ that are computed in the $(k-1)$-th round. Moreover, the contribution of any $w_j^{(k-1)}$ with $j\neq i$ to $u_i^{(k)}$ is proportional to the similarity between $w_i^{(k-1)}$ and $w_j^{(k-1)}$. Before ending this subsection, we provide in Algorithm \ref{alg:Fed-GK} the pseudo-code of {FedAMP} with $D$ being the Gaussian kernel distance function. 

% \begin{algorithm}[t]
% \DontPrintSemicolon
% \SetKwInput{KwInput}{Input}               
% \SetKwInput{KwOutput}{Output} 
% \KwInput{number of clients $m$, initial models $(w_1^{(0)},\dots, w_m^{(0)})$, total communication rounds $K$, self-attention parameters $\{\tau_{k,i}\}\subset(0,1]$, local training parameters $\{\beta_{k,i}\}\subset\R_{++}$}
% \KwOutput{the final models $(w_1^{(K)},\dots, w_m^{(K)})$ and $(u_1^{(K)},\dots, u_m^{(K)})$}             
% \For{$k=1,2,\dots, K$}{
% for every $i=1,\dots,m$, server computes the attentive aggregation of $\{w_j^{(k-1)}\}_{j\neq i}$ by 
% \begin{equation}\label{eq:aggre}
% w^{(k-1)}_{i,c} = \frac{\sum_{j\neq i}^m e^{-\frac{\|w_i^{(k-1)} - w_j^{(k-1)}\|^2}{2\sigma}}\cdot w_j^{(k-1)}}{\sum_{j\neq i}^m e^{-\frac{\|w_i^{(k-1)} - w_j^{(k-1)}\|^2}{2\sigma}}}
% \end{equation}\\
% for every $i=1,\dots,m$, servers computes the intermediate model $u_i^{(k)}$ by
% \begin{equation}\label{eq:cvx-comb}
% u_i^{(k)} = (1 - \tau_{k,i})w_i^{(k-1)} + \tau_{k,i}z_i^{(k-1)}
% \end{equation}\\
% server sends each $u_i^{(k)}$ to client $i$ \\
% clients perform local training and compute the updated models $(w_1^{(k)},\dots,w_m^{(k)})$ by \eqref{eq:local-training}\\
% %\begin{equation}
% %\label{eq:local-training-c}
% %w_i^{(k)} \approx \arg\min_{w\in\R^d} F_i(w) + \frac{1}{2\alpha_{k,i}}\|w - u_i^{(k)}\|^2
% %\end{equation}
% server collects the updated models $(w_1^{(k)},\dots,w_m^{(k)})$
% }
% \caption{{FedAMP} with Gaussian Kernel Distance Function}
% \label{alg:Fed-GK}
% \end{algorithm}

\nop{
conduct privacy-preserving machine learning under the coordination of a service provider
}